% titel: Blatt 8, A7: Hydrogensalze besser löslich als Normalsalze
% skript: Mortimer Kap. 12 (Löslichkeit, Hydratationsenergie)

% Teil I: Werkzeuge, kein Lösungsweg.

Warum löst sich ein Hydrogensalz wie NaHSO$_4$ (oder NaHCO$_3$) meist besser
in Wasser als das zugehörige Normalsalz? Das hängt mit der Energiebilanz des
Lösungsvorgangs zusammen.

\begin{werkzeugkasten}[Löslichkeit und Energiebilanz]
Lösen eines Salzgitters in Ionen:
\[
  \Delta H_{\text{Lösung}} = \Delta H_{\text{Gitter}}
  - \Delta H_{\text{Hydratation}} .
\]
Je \emph{kleiner} die Gitterenergie (leicht aufzubrechen) und je
\emph{größer} die Hydratationsenergie (Wasser-Ionen-Anziehung), desto
leichter löst sich das Salz.

Das Anion HSO$_4^-$ (oder HCO$_3^-$):
\begin{itemize}
  \item ist \emph{einfach geladen}: weniger Gitterenergie als das
        zweifach geladene SO$_4^{2-}$,
  \item trägt ein H-Atom, das \textbf{Wasserstoffbrücken} zum Wasser
        ausbilden kann $\Rightarrow$ hohe Hydratation.
\end{itemize}
\textit{Skript:} Mortimer Kap. 12 (Löslichkeit von Salzen).
\end{werkzeugkasten}

\begin{achtung}[Darauf achten]
\begin{itemize}
  \item Vergleiche immer die \emph{Gitterenergie}: zweifach geladene
        Anionen (SO$_4^{2-}$) haben eine viel größere Coulomb-anziehung als
        einzach geladene (HSO$_4^-$).
  \item Der Wassersstoff der Hydrogen-Anionen ermöglicht zusätzliche
        H-Brücken zum Lösungsmittel.
  \item Es gibt Ausnahmen --- die Aussage im Blatt nennt ``abgesehen von
        einigen Ausnahmen``, man beschreibt also den allgemeinen Trend.
\end{itemize}
\end{achtung}
